% Propietario conservado: /Users/ruben/Documents/ChatGPT/jueces y controles/REVISION_ARGUMENTAL_APERTURAS_CIERRES_20260915/delivery/ARTICLE_V_ES_ARGUMENTAL_20260915/manuscrito/sections/30_fock_gibbs.tex
% Recuperación documental para el artículo V; RESULTADO_RECUPERADO.
% Propietario: /Users/ruben/Documents/New project/output/REVISION_EDITORIAL_HMT_MD_20260905/01_FUENTES/INTEGRAL/tree/New project/output/ENSAMBLADO_CAUSAL_RESTAURADO_HMT_MD_20260905/fuente/manuscrito/sucesor_102/deltas_masas/realizacion_equivariante_especies_operadores_masicos_publico.tex
% Líneas de origen: 856--932.
% REV02: transporte functorial; CAR/CCR se construyen en 20.
% Correspondencia editorial: gestion/CONSOLIDACION_PAULI_FOCK_REV02.json.
\section{Transport des complétions de Fock par raffinement}
\label{v:v:rec:fock_gibbs:section:01}

Les complétions et leurs algèbres d’opérateurs sont construites dans la
section~\ref{v:v:sec:completaciones-algebraicas}. Le résultat de cette étape
est leur compatibilité avec les morphismes de raffinement : le transport de
la fibre à une particule se prolonge à tous les secteurs d’occupation.

\subsection{Seconde quantification des morphismes de raffinement}
\label{v:v:rec:fock_gibbs:section:02}

Soit \(\mathsf{Hol}_{\pm}^{\leq1}\) la catégorie monoïdale des fibres à une
particule, avec des morphismes de raffinement contractifs qui conservent la
parité d’holonomie \(\epsilon\in\{+1,-1\}\). Elle inclut, en particulier,
les entrelaceurs unitaires. L’opération monoïdale sur les fibres est
la somme directe ; dans les complétions, elle induit le produit tensoriel,
gradué dans le secteur extérieur. La seconde quantification utilise les
complétions de la section~\ref{v:v:rec:completaciones:section:03} :
\[
 \mathfrak F(H,\epsilon)=
 \begin{cases}
 \mathcal F_+(H),
     &\epsilon=+1,\\[2mm]
 \mathcal F_-(H),
     &\epsilon=-1,
 \end{cases}
\]
et, pour un entrelaceur contractif \(T:H\to H'\),
\[
 \mathfrak F(T)=
 \bigoplus_{n\ge0}\operatorname{Sym}^n(T)
 \quad\text{ou}\quad
 \bigoplus_{n\ge0}\bigwedge^n\!T.
\]

\subsection{Fonctorialité, opérateurs canoniques et domaines de transport}
\label{v:v:rec:fock_gibbs:section:03}

Soit \(\mathcal D_\epsilon(H)\) la somme directe algébrique des secteurs
à nombre fini de particules. Dans les identités suivantes, \(a_H^\dagger(f)\)
et \(a_H(f)\) désignent les opérateurs de création et d’annihilation de la
complétion choisie : les opérateurs \(a\) du secteur extérieur ou les
opérateurs \(b\) du secteur symétrique de la
section~\ref{v:v:pauli:car-construccion}.

\begin{theorem}[Transport interne de la statistique]
\label{v:v:rec:fock_gibbs:statement:01}
La complétion précédente est un foncteur d’opérateurs bornés sur
\(\mathsf{Hol}_{\pm}^{\leq1}\), compatible avec le raffinement et
monoïdal gradué. Dans la branche \(\epsilon=-1\), les opérateurs de création et
d’annihilation satisfont les CAR et les opérateurs d’occupation satisfont
\(N_s^2=N_s\) ; dans la branche \(\epsilon=+1\), ils satisfont les CCR. Par conséquent,
l’exclusion de Pauli et les statistiques fermionique et bosonique descendent de la
parité d’holonomie une fois fixée la représentation à une particule.
Pour tout entrelaceur contractif \(T:H\to H'\), \(f\in H\) et
\(g\in H'\), on a dans \(\mathcal D_\epsilon(H)\)
\begin{align}
 \mathfrak F(T)a_H^\dagger(f)
 &=a_{H'}^\dagger(Tf)\mathfrak F(T),
 \label{v:v:fock:eq:transporte-creacion}\\
 a_{H'}(g)\mathfrak F(T)
 &=\mathfrak F(T)a_H(T^*g).
 \label{v:v:fock:eq:transporte-aniquilacion}
\end{align}
Si \(T\) est isométrique, la seconde identité, avec \(g=Tf\), transporte
aussi l’annihilation de \(f\). Pour \(\|f\|=1\), les opérateurs
d’occupation satisfont alors
\begin{equation}
 N_{H',Tf}\mathfrak F(T)=\mathfrak F(T)N_{H,f},
 \qquad N_{H,f}=a_H^\dagger(f)a_H(f).
 \label{v:v:fock:eq:transporte-ocupacion}
\end{equation}
\end{theorem}

\begin{proof}
Les puissances symétriques et extérieures préservent les compositions et les identités ;
à chaque degré, leurs normes sont bornées par \(\|T\|^n\). La contractivité
garantit ainsi que la somme directe définit un opérateur borné. La
décomposition des puissances d’une somme directe selon leurs degrés dans chaque
facteur donne la compatibilité monoïdale, avec le signe d’échange gradué
dans les puissances extérieures.

Les relations CAR et l’idempotence sont le
théorème~\ref{v:v:rec:completaciones:statement:01} ; les relations CCR sont
\eqref{v:v:pauli:eq:ccr}. Pour vérifier leur transport, il suffit d’agir sur
des produits symétriques ou extérieurs de vecteurs. Appliquer \(T\) à chaque facteur
commute avec l’insertion du facteur \(f\), qui se transforme en \(Tf\) ;
cela prouve \eqref{v:v:fock:eq:transporte-creacion}. L’annihilation somme les
contractions de chaque facteur et utilise
\(\langle g,Tf\rangle=\langle T^*g,f\rangle\) ; les signes extérieurs et
les normalisations symétriques coïncident des deux côtés. On obtient
\eqref{v:v:fock:eq:transporte-aniquilacion}. Si \(T^*T=I\), poser \(g=Tf\)
et composer les deux identités démontre
\eqref{v:v:fock:eq:transporte-ocupacion}. Toutes ces opérations préservent
le domaine à nombre fini de particules.
\end{proof}

Pour les morphismes non contractifs, la seconde quantification est d’abord définie
dans \(\mathcal D_\epsilon(H)\) ; son extension bornée dépend de la borne
uniforme des normes de tous les degrés. De tels morphismes restent hors
de \(\mathsf{Hol}_{\pm}^{\leq1}\). Dans la complétion symétrique, si
\(\|T\|>1\), les normes \(\|\operatorname{Sym}^nT\|=\|T\|^n\)
croissent sans borne ; dans la complétion extérieure d’une fibre finie, il n’y a
qu’un nombre fini de degrés non nuls. L’identité générale
d’annihilation conserve \(T^*g\) ; sa spécialisation au même vecteur
transporté correspond à la condition isométrique indiquée.

La distinction se vérifie déjà dans une fibre unidimensionnelle. Si
\(T=tI\), avec \(0<t<1\), et que \(e\) est unitaire, alors, sur l’état
à une particule \(e\), dans l’une ou l’autre des deux complétions,
\begin{equation}
 a(te)\mathfrak F(T)e=t^2\Omega,
 \qquad \mathfrak F(T)a(e)e=\Omega,
 \label{v:v:fock:eq:control-no-isometrico}
\end{equation}
où \(\Omega\) est le vide. L’identité générale
\eqref{v:v:fock:eq:transporte-aniquilacion} donne correctement
\(\mathfrak F(T)a(T^*te)e=t^2\Omega\). Le facteur \(t^2\) identifie
exactement l’hypothèse qui permet de spécialiser le transport de
l’annihilation puis celui de l’occupation.

Le résultat établit le transport algébrique interne. Son identification
au théorème relativiste de spin--statistique conserve les conditions
supplémentaires de représentation lorentzienne et de localité explicitées dans la
section~\ref{v:v:rec:completaciones:section:05}.

\clearpage
\section{Réalisation thermodynamique au moyen d’états de Gibbs et limite de pression}
\label{v:v:rec:fock_gibbs:section:04}

\subsection{Hamiltonien fini, partition, état de Gibbs et limite cofinale de la pression}
\label{v:v:rec:fock_gibbs:section:05}

Pour un Hamiltonien comprimé \(H_{X,V}\) sur un volume ou un raffinement
fini \(V\), l’état thermique est complètement défini par
\[
 Z_{X,V}(\beta)=\operatorname{Tr}e^{-\beta H_{X,V}},
 \qquad
 \rho_{X,V,\beta}=Z_{X,V}(\beta)^{-1}e^{-\beta H_{X,V}}.
\]
La limite thermodynamique existe lorsque la pression
\[
 p_X(\beta)=\lim_{V\nearrow\infty}
 \frac1{\beta |V|}\log Z_{X,V}(\beta)
\]
existe et est indépendante de la suite cofinale.

\subsection{Critère de condensation par saturation de la densité excitée}
\label{v:v:rec:fock_gibbs:section:06}

La condensation bosonique est
décidée par la saturation de la densité des états excités ; elle n’est pas proclamée à
partir de la seule existence d’une complétion symétrique. Ainsi, le problème
est exprimé au moyen d’une condition mathématique vérifiable pour le
Hamiltonien et la limite, au lieu d’une attribution nominale.
